\section{Introduction}
Atopic dermatitis treatment must improve visible skin disease, itch and the burden of living with eczema. Treat-to-target frameworks bring these complementary outcomes into treatment decisions, including early symptom relief and later disease control.\citep{williams2022home,debruin2021t2t,zhao2026t2t} A practical severity estimate should remain interpretable as treatment changes and should capture how an individual patient's condition changes, rather than only distinguish groups with different average severity.

Serum markers offer a link between clinical activity and inflammation. The p-EASI approach combines CCL17, IL-22 and soluble IL-2 receptor, with different coefficients before and during treatment.\citep{thijs2017} Subsequent studies reported Spearman correlations of 0.71 during cyclosporin treatment and 0.67 during dupilumab treatment.\citep{thijs2019,bakker2020} These studies established a practical precedent for converting blood measurements into Eczema Area and Severity Index (EASI) units. Applying such a calculation in another setting also requires agreement with individual scores, compatible assays and evaluation across the severity range.\citep{luijken2020,riley2024validation}

EASI combines the extent of affected skin with lesion intensity.\citep{hanifin2022,home_easi} A small-input equation combining affected body surface area (BSA) and serum measurements could offer a common calculation without separate rules for each drug or treatment phase. BSA supplies the clinical extent component; the serum panel supplies molecular information. We evaluated it as a research readout of recorded EASI at the visit, with separate tests of individual change and molecular added accuracy. Measured itch, quality of life and safety remained complementary clinical assessments.

We investigated these questions in a Chinese cohort receiving dupilumab or two initial abrocitinib doses. We first established the clinical and molecular context, retaining all 23 measured serum markers and accounting for baseline differences and later treatment changes. We then developed a common BSA--serum equation, evaluated individual estimates and changes using patient-level validation, and examined unchanged published p-EASI formulas as a separate local-application comparison.
